\chapter{Build a Research Workflow}\label{ch:rs:build-a-research-workflow}

A result from eighteen months ago cannot be reproduced. The vendor has revised its history since, the code has moved on, and nobody kept the random seed of the
bootstrap that produced the confidence interval in the committee's slides. Each of the three is ordinary, and together they make the result unknowable: it may
have been right, but nobody can say. This last chapter of the book puts its reference study, from a \omterm{def:rs:build-a-research-workflow:repro}{data snapshot} to a tear sheet, into a pipeline whose every
stage is keyed by the hash of its code, parameters, seed and inputs, so that a result is a manifest that can be re-run bit for bit, and a change anywhere says
exactly what it changed. The build is \texttt{firm.workflow}.

\section{Reproducibility}

\begin{definition}[Reproducible result, data snapshot, content-addressed storage, pipeline stage, environment lock]\label{def:rs:build-a-research-workflow:repro}
A \emph{reproducible result}\index{reproducible result} is one that the same data, code and settings produce again, bit for bit (Book~4, chapter~25), a minimum
standard when an independent replication is not possible (Peng). A \emph{data snapshot}\index{data snapshot} is an immutable copy of the data as used, identified
by its content. \emph{Content-addressed storage}\index{content-addressed storage} stores each object under the hash of its content, so that the same content has the
same address and different content cannot share one (Git's object store is the familiar example). A \emph{pipeline stage}\index{pipeline stage} is a function of its
inputs, parameters and seed, with no other source of variation. An \emph{environment lock}\index{environment lock} records the exact versions of the interpreter
and libraries a result was produced with.
\end{definition}

Sandve and co-authors reduced the practice to ten rules; the first, sixth and ninth carry most of it: for every result, keep track of how it was produced; for
analyses that include randomness, note the seeds; connect textual statements to the results behind them. The book has applied the ninth rule to itself since
chapter~1: every number printed in these chapters is asserted by a test that recomputes it. What a research desk adds is scale: many studies, shared stages,
vendor data that changes under them. The answer is to make the pipeline, not the analyst, remember how each result was produced.

\texttt{firm.workflow} keys each stage by the SHA-256 of its name, the hash of its source code, its parameters, its seed and the content hashes of its inputs'
outputs, stores the output under that key, and hashes the output in turn (\cref{lst:rs:build-a-research-workflow:run}). A stage whose key is already stored does
not run. Because a key depends on its inputs' outputs rather than their keys, a change upstream that leaves an intermediate result unchanged stops there. A run
writes a manifest: for each stage its key, code hash, parameters, seed, input hashes, output hash and whether it ran, and the environment.

\section{Review}

\begin{definition}[Review checklist]\label{def:rs:build-a-research-workflow:review}
A \emph{review checklist}\index{review checklist} is the list of questions a \omterm{def:rs:the-research-process:gate}{research review} (chapter~1) answers before a result is accepted: that it is
reproduced from its manifest, that its data are point in time, that its trials are counted, that its \omterm{def:rs:vectorised-backtests:fidelity}{fidelity level} matches the claim, and that its costs and
\omterm{def:rs:capacity-decay-and-crowding:capacity}{capacity} are those of the size it will run at.
\end{definition}

The book's chapters are its checklist, and a reviewer can ask their questions of any manifest.
\begin{itemize}
\item \textbf{Reproduced.} Does the manifest reproduce bit for bit in an empty cache, in the locked environment? Which stages have no seed?
\item \textbf{Point in time.} Is the snapshot identified, and are its revisions recorded (chapter~3)? Are the universe and the exposures known at each date
(chapters~4 and~24)?
\item \textbf{Counted.} How many trials are in the \omterm{def:rs:the-research-process:log}{research log}'s family, and what are the deflated Sharpe ratio and the \omterm{def:rs:overfitting:pbo}{probability of backtest overfitting}
(chapters~1 and~20)?
\item \textbf{Faithful.} At which \omterm{def:rs:vectorised-backtests:fidelity}{fidelity level} was it backtested, and is the claim within what that level can say (chapters~16 to~19)?
\item \textbf{Costed and sized.} Are the costs those of the fitted model at the intended size, and where is the \omterm{def:rs:capacity-decay-and-crowding:capacity}{capacity} (chapters~27 and~28)?
\item \textbf{Stated.} Is every number in the write-up traceable to a stage's output, and every chart to its data (Sandve's ninth and seventh rules)?
\end{itemize}
Arnott, Harvey and Markowitz set out a similar protocol for research with machine learning; the pipeline makes its first questions mechanical, so that the review
can spend its time on the last ones.

\section{Build: the pipeline from raw data to a reviewed result}

The reference study has eight stages (\cref{fig:rs:build-a-research-workflow:dag}, \cref{lst:rs:build-a-research-workflow:make}): the snapshot of
\texttt{firm.synthmkt} with any revisions the vendor has made; the 500 most liquid names each day; two features, 12--1 month momentum and the one-day reversal;
their predictor-card statistics; an equal blend; a monthly dollar-neutral book; a level-1 \omterm{def:rs:vectorised-backtests:backtest}{backtest} with costs of ten basis points; and a tear sheet with a block
bootstrap of the Sharpe ratio, seeded. The first run computes every stage in a few seconds; the second reads everything from the cache.

\begin{omfigure}
\begin{tikzpicture}[st/.style={draw=omDef,fill=omDef!8,rounded corners=2pt,font=\scriptsize,minimum width=1.7cm,minimum height=0.7cm,align=center},
  hit/.style={st,draw=omThm,fill=omThm!12},arr/.style={-{Stealth},thick,black!60}]
\node[hit] (s) at (0,0) {snapshot};
\node[hit] (u) at (2.1,0) {universe};
\node[hit] (f) at (4.2,0) {features};
\node[hit] (c) at (6.3,1.3) {cards};
\node[st] (b) at (6.3,0) {blend};
\node[st] (p) at (8.4,0) {portfolio};
\node[hit] (l) at (10.5,0) {level 1};
\node[st] (t) at (12.6,0) {tear sheet};
\foreach \a/\b in {s/u,u/f,f/c,f/b,b/p,p/l,l/t} \draw[arr] (\a) -- (\b);
\draw[arr] (s.north) to[out=50,in=130] (f.north);
\draw[arr] (s.north east) to[out=70,in=180] (c.west);
\draw[arr] (u.south) to[out=-50,in=-130] (b.south);
\draw[arr] (u.south east) to[out=-60,in=-120] (p.south);
\draw[arr] (s.south) to[out=-40,in=-140,looseness=0.6] (l.south);
\end{tikzpicture}

\omcaption{The reference study's eight stages. A vendor revision to a name that never enters the universe re-runs the five red stages, but only the snapshot's
output changes: the other four reproduce their previous outputs, and blend, portfolio and tear sheet are read from the cache.}\label{fig:rs:build-a-research-workflow:dag}
\end{omfigure}

The study's own result is modest, and the pipeline reports it with its uncertainty: momentum's mean daily \omterm{def:rs:anatomy-of-a-predictor:ic}{rank information coefficient} is 0.0076 ($t = 1.39$), the
reversal's 0.0379 ($t = 16.05$) over 2\,267 days, and the blended monthly book has a Sharpe ratio after costs of 0.08 (standard error 0.33; bootstrap interval from
$-0.67$ to 0.74), turning over 3.5\% of its capital a day. The reversal's strong daily signal does not survive a book rebuilt monthly, the horizon mismatch of
chapter~6. What matters here is how the result behaves under change:

\begin{center}
\small
\begin{tabular}{@{}lll@{}}
\toprule
change & stages re-run & outputs changed\\
\midrule
none & none & none\\
revision, name never in the universe & snapshot, universe, features, cards, level 1 & snapshot\\
revision, name always in the universe & all eight & all eight\\
blend weights 0.7 and 0.3 & blend, portfolio, level 1, tear sheet & those four\\
tear sheet code gains a \omterm{def:rs:performance-measurement:hit}{hit rate} & tear sheet & tear sheet\\
\bottomrule
\end{tabular}
\end{center}

A revision to a name that is never among the liquid 500 re-runs every stage that reads the snapshot, and each of them reproduces its previous output, so nothing
downstream of them runs: the manifest diff shows the snapshot's parameters and output changed and the others' inputs only. A revision inside the universe changes
everything, and says so. A change of code changes only its stage. And reproduction, re-running the whole manifest in an empty cache, returns the same output hash
for all eight stages; with the bootstrap's seed left out it flags exactly one, the tear sheet (\cref{lst:rs:build-a-research-workflow:repro}). Finally the run is
registered as a trial in the \omterm{def:rs:the-research-process:log}{research log} of chapter~1, with the hash of its manifest as its code identity and the snapshot's hash as its data identity: the
eighteen-month-old result, had it been produced this way, would be one line of the log and one command.

\section{Tutorial: the result that could not be reproduced}\label{tut:rs:build-a-research-workflow:tut}

\begin{tutorial}
\textbf{Goal.} Run the reference study through \texttt{firm.workflow}, change the data, the parameters and the code in turn, reproduce the result bit for bit, and
register it. \textbf{End state:} the table and \cref{fig:rs:build-a-research-workflow:dag}.
\begin{tutsteps}
\item \textbf{Running a pipeline}: keys from code, parameters, seed and input hashes; the cache; the manifest.
\omcode{code/firm/workflow/firm_workflow.py}{107}{123}{A content-addressed pipeline run.}\label{lst:rs:build-a-research-workflow:run}
\item \textbf{Reproduction and diff.}
\omcode{code/firm/workflow/firm_workflow.py}{126}{142}{Re-running from scratch, and what changed between two manifests.}\label{lst:rs:build-a-research-workflow:repro}
\item \textbf{The reference study} as eight stages.
\omcode{code/research/29-build-a-research-workflow/python/rs_workflow.py}{115}{129}{The study's pipeline.}\label{lst:rs:build-a-research-workflow:make}
\item \textbf{Run} \texttt{rs\_workflow.make()} in a cache directory; revise a name with \texttt{never\_in\_universe} and \texttt{always\_in\_universe}; change the
weights; swap in \texttt{tearsheet\_v2}; call \texttt{reproduce} and \texttt{register}.
\end{tutsteps}
\textbf{What to change next.} Add a level-2 stage (chapter~17) on the book's largest name and see which changes reach it; store the cache in a shared directory and
let a colleague reproduce your manifest on another machine, comparing environments first.
\end{tutorial}

\section{Build: the workflow}\label{bld:rs:build-a-research-workflow:workflow}

\begin{build}
\bfield{Purpose} Every result the firm relies on is produced by a pipeline whose manifest reproduces it, says what it depends on, and is registered in the \omterm{def:rs:the-research-process:log}{research
log}.
\bfield{Interface} \texttt{content\_hash(obj)}, \texttt{Stage(name, fn, deps, params, seed)}, \texttt{Pipeline(stages, cache\_dir).run(target)} returning outputs and
a manifest, \texttt{environment()}, \texttt{reproduce(make\_pipeline, manifest, cache\_dir)}, \texttt{diff(m1, m2)}, \texttt{register(manifest, log, family, ts,
metrics)}.
\bfield{Rules} Stages are pure functions of inputs, parameters and seed; data enter only through snapshot stages whose revisions are parameters; every stage that
draws random numbers takes a seed; manifests and caches are kept with the result; a result is reviewed from its manifest, not from its slides.
\bfield{Acceptance tests} \texttt{code/firm/workflow/tests/}: a canonical hash (key order irrelevant, types significant); a second run entirely from the cache; early
cutoff when a revision is clipped away; diffs naming the changed fields; reproduction bit for bit, and a failed reproduction flagging exactly the unseeded stage;
registration in the \omterm{def:rs:the-research-process:log}{research log}.
\bfield{Stretch} Code hashes that follow the functions a stage calls; remote caches shared by the desk; \omterm{def:rs:build-a-research-workflow:repro}{environment locks} enforced before reproduction.
\end{build}

\begin{omsources}
\item R.~D.~Peng, ``Reproducible research in computational science'', \emph{Science} 334(6060), 2011.
\item G.~K.~Sandve, A.~Nekrutenko, J.~Taylor and E.~Hovig, ``Ten simple rules for reproducible computational research'', \emph{PLoS Computational Biology} 9(10),
2013.
\item \emph{Pro Git}, 2nd edition, section 10.2, ``Git internals: Git objects'' (git-scm.com).
\item R.~Arnott, C.~R.~Harvey and H.~Markowitz, ``A backtesting protocol in the era of machine learning'', \emph{Journal of Financial Data Science} 1(1), 2019.
\end{omsources}

\section{Exercises}

\begin{exercise}[$\star$]\label{exo:rs:build-a-research-workflow:1}
Why must a stage's key include the hash of its code, and what does a key that includes only its parameters miss?
\end{exercise}

\begin{exercise}[$\star$]\label{exo:rs:build-a-research-workflow:2}
Why are keys built from the inputs' output hashes rather than the inputs' keys?
\end{exercise}

\begin{exercise}[$\star$]\label{exo:rs:build-a-research-workflow:3}
The vendor revises a return of a stock that was never in the universe. Which stages run, and which outputs change?
\end{exercise}

\begin{exercise}[$\star\star$]\label{exo:rs:build-a-research-workflow:4}
Which of Sandve's ten rules does the pipeline enforce, and which does it leave to people?
\end{exercise}

\begin{exercise}[$\star\star$]\label{exo:rs:build-a-research-workflow:5}
A stage calls a helper function in another module. Its code hash is the hash of the stage's own source. What can go wrong, and how would you fix it?
\end{exercise}

\begin{exercise}[$\star\star$]\label{exo:rs:build-a-research-workflow:6}
Reproduction fails on a colleague's machine for every stage that uses floating-point reductions. What do you check first?
\end{exercise}

\begin{exercise}[$\star\star\star$]\label{exo:rs:build-a-research-workflow:7}
\emph{Coding.} Remove the tear sheet's seed and run \texttt{reproduce}. Then give the bootstrap a seed derived from the manifest's other keys and explain why that is
reproducible but still honest.
\end{exercise}

\begin{exercise}[$\star\star\star$]\label{exo:rs:build-a-research-workflow:8}
\emph{Find the flaw.} ``The study is reproducible: the notebook is in version control.''
\end{exercise}

\section{Problem: The Result That Could Not Be Reproduced}

\begin{problem}[{Weekend problem --- a result, kept}]\label{pb:rs:build-a-research-workflow:1}
The reference study in \texttt{firm.workflow}.

\medskip\noindent\textbf{Part I --- The pipeline.}
\begin{enumerate}
\item List the eight stages and their dependencies.
\item What goes into a stage's key, and what into the manifest?
\item What does the second run do, and why?
\item What does the study find, with its uncertainty?
\end{enumerate}

\medskip\noindent\textbf{Part II --- Changes.}
\begin{enumerate}[resume]
\item Which stages re-run after a revision to a name never in the universe, and which outputs change?
\item And after a revision to a name always in the universe?
\item After a change of blend weights?
\item After a change to the tear sheet's code?
\end{enumerate}

\medskip\noindent\textbf{Part III --- Reproduction.}
\begin{enumerate}[resume]
\item What does reproducing the manifest in an empty cache return?
\item What happens without the bootstrap's seed?
\item What does the environment record, and when does it matter?
\end{enumerate}

\medskip\noindent\textbf{Part IV --- The verdict.}
\begin{enumerate}[resume]
\item State the \emph{named result}: the stages invalidated by a vendor revision, and the manifest diff that proves the reproduction.
\item How is the run registered in the \omterm{def:rs:the-research-process:log}{research log}?
\item Which items of the \omterm{def:rs:build-a-research-workflow:review}{review checklist} does the pipeline answer, and which not?
\item Why does the reversal's daily signal not survive the monthly book?
\item How would the eighteen-month-old result have been kept?
\item What would you add before the desk relies on the pipeline?
\item In one sentence: what is a \omterm{def:rs:build-a-research-workflow:repro}{reproducible result}?
\end{enumerate}
\end{problem}

\section{Interview questions}

\begin{interviewq}[$\star$ \iqroles{researcher, developer}]\label{iq:rs:build-a-research-workflow:1}
A \omterm{def:rs:vectorised-backtests:backtest}{backtest} you ran last year gives a different Sharpe ratio today. List the possible causes.
\end{interviewq}

\begin{interviewq}[$\star\star$ \iqroles{researcher}]\label{iq:rs:build-a-research-workflow:2}
What practices make computational research reproducible?
\end{interviewq}

\begin{interviewq}[$\star\star$ \iqroles{developer}]\label{iq:rs:build-a-research-workflow:3}
Design a cache for research \omterm{def:rs:build-a-research-workflow:repro}{pipeline stages}. How do you invalidate it?
\end{interviewq}

\begin{interviewq}[$\star\star$ \iqroles{researcher, risk}]\label{iq:rs:build-a-research-workflow:4}
What would you check when reviewing a colleague's strategy research?
\end{interviewq}

\begin{interviewq}[$\star\star$ \iqroles{developer}]\label{iq:rs:build-a-research-workflow:5}
How do you make a computation that uses random numbers reproducible across machines?
\end{interviewq}

\begin{interviewq}[$\star\star\star$ \iqroles{developer, researcher}]\label{iq:rs:build-a-research-workflow:6}
Your vendor restates five years of history. How do you find every result that depended on it, and which of them change?
\end{interviewq}
